\documentclass[11pt,a4paper]{article}
\usepackage[T1]{fontenc}
\usepackage{lmodern,amsmath,amssymb,amsthm,mathtools,booktabs,microtype}
\usepackage[margin=26mm]{geometry}
\usepackage[colorlinks=true,linkcolor=blue,citecolor=blue,urlcolor=blue]{hyperref}
\setlength{\emergencystretch}{2em}
\newtheorem{theorem}{Theorem}[section]
\newtheorem{proposition}[theorem]{Proposition}
\newtheorem{lemma}[theorem]{Lemma}
\newtheorem{corollary}[theorem]{Corollary}
\theoremstyle{definition}\newtheorem{definition}[theorem]{Definition}
\theoremstyle{remark}\newtheorem{remark}[theorem]{Remark}
\newcommand{\Q}{\mathbb Q}\newcommand{\Z}{\mathbb Z}
\newcommand{\R}{\mathbb R}\newcommand{\A}{\mathbb A}
\newcommand{\PP}{\mathbb P}\newcommand{\sqcl}[1]{[#1]}
\DeclareMathOperator{\Tr}{Tr}\DeclareMathOperator{\Nm}{N}
\DeclareMathOperator{\N}{N}\DeclareMathOperator{\rad}{rad}
\DeclareMathOperator{\sfpart}{sf}
\title{Unchanged Existential Formulas after Finite Field Extensions}
\author{}\date{16 November 2026}
\begin{document}\maketitle
\begin{abstract}
An existential formula with infinitely many inputs in a number field
accepts every input outside a finite exceptional set after a field
extension of uniformly bounded degree. Consequently no unchanged
formula defines the rings of integers, or the corresponding rings of
$S$-integers, in all finite extensions. A fixed finite collection of
formulas fails simultaneously in extensions with a uniform Galois
degree bound. For one nonzero polynomial, the required extension
degree is at most its smallest positive individual witness degree.
An exact quadratic-extension specialization illustrates the obstruction.
These results concern reusing fixed formulas and do not exclude
field-dependent existential definitions.
\end{abstract}

\section{Existential formulas after finite extensions}

Fix a number field $k$. An existential formula considered here
has the form
\[
 \Phi(T):\quad \exists Y_1,\ldots,Y_m\quad
             \theta(T,Y_1,\ldots,Y_m),
\]
where $\theta$ is a finite Boolean combination of polynomial
equations with coefficients in $k$. Equivalently, it is a finite
disjunction of conjunctions of equations and nonzero conditions.
For each extension field $L/k$, write
\[
                \Phi(L)=\{a\in L:\Phi(a)\text{ holds in }L\}.
\]
The same formula and the same coefficients are used for every $L$.

\begin{theorem}\label{sys:bounded-extension}
If $\Phi(k)$ is infinite, then there exist a positive integer
$D$ and a nonzero polynomial $h\in k[T]$ such that every
$a\in k$ with $h(a)\ne0$ belongs to $\Phi(L_a)$ for some finite
extension $L_a/k$ of degree at most $D$.
\end{theorem}
\begin{proof}
The conditions defining the witness tuples give a constructible
subset of affine space. By Chevalley's theorem its image in
the input line is constructible \cite{StacksChevalley}. Since
the image has infinitely many $k$-points, it contains the
generic point of the line. Equivalently, at least one of the
finitely many conjunctions has a nonempty generic fiber over
$k(T)$.

Choose a closed point in a nonempty affine open of that generic
fiber. Its residue field $E$ is a finite extension of $k(T)$.
The coordinates of the point give elements
$y_1,\ldots,y_m\in E$ satisfying one conjunction, with every
required nonzero expression nonzero in $E$. Characteristic zero
gives a primitive element $\alpha$ for $E/k(T)$. Let
$G(T,Z)\in k(T)[Z]$ be its monic minimal polynomial, of
degree $D=[E:k(T)]$ in $Z$.

Express each $y_i$, and also the inverse of each required
nonzero expression, as a polynomial in $\alpha$ of degree
less than $D$ with coefficients in $k(T)$. Choose a nonzero
$h\in k[T]$ which clears all denominators in these expressions
and in $G$. Put $R=k[T,1/h]$ and
\[
                              C=R[Z]/(G).
\]
Because $G$ is monic, $C$ is free of rank $D$ over $R$ and
embeds in $E$ under $Z\mapsto\alpha$. The chosen coordinate
expressions and inverse expressions belong to $C$. All the
equations, and the identities asserting that the chosen inverses
are inverses, hold in $C$: they hold in $E$, and $C\hookrightarrow E$
is injective.

Now let $a\in k$ with $h(a)\ne0$. Specialization gives the
nonzero $D$-dimensional algebra
\[
                             C_a=k[Z]/(G(a,Z)).
\]
Choose a maximal ideal of $C_a$. Its residue field $L_a$ is a
finite extension of $k$ of degree at most $D$. The coordinate
expressions descend to a witness tuple in $L_a$; all equations
and inverse identities survive. Hence the same conjunction,
and therefore $\Phi(a)$, holds in $L_a$.
\end{proof}

The finite exceptional set in the theorem consists of the
$k$-roots of the single fixed polynomial $h$. The degree bound
$D$ is uniform in $a$, although the extension $L_a$ may depend
on $a$. The proof specifies $D$ as the degree of one generic
closed point; it does not identify one extension that contains
witnesses for all inputs.

\begin{corollary}\label{sys:no-uniform-integers}
No single existential formula over a number field $k$ defines
$\mathcal O_L$ in every finite extension $L/k$. More precisely,
if $\Phi(k)$ is infinite, there is a positive integer $D$ such
that, for all but finitely many $q\in\mathbb Q\setminus\mathbb Z$,
the formula accepts $q$ in some extension $L_q/k$ of degree at
most $D$.
\end{corollary}
\begin{proof}
Apply Theorem~\ref{sys:bounded-extension} to the elements
$q\in\mathbb Q\subset k$. Every rational algebraic integer is
an integer. Indeed, if $q=a/b$ is in lowest terms and satisfies
a monic equation of degree $r$ with integer coefficients,
clearing denominators yields $b\mid a^r$, so $b=1$.
Consequently each admitted rational noninteger lies outside
$\mathcal O_{L_q}$.
\end{proof}

\begin{corollary}\label{sys:no-uniform-s-integers}
Let $S$ be a fixed finite set of finite primes of $k$, and for
each finite extension $L/k$ let $S_L$ be the primes of $L$
above $S$. One unchanged existential formula over $k$ cannot
define $\mathcal O_{L,S_L}$ in every finite extension $L/k$.
\end{corollary}
\begin{proof}
Such a formula would have infinitely many inputs over $k$,
since $\mathbb Z\subseteq\mathcal O_{k,S}$. Choose a rational
prime $p$ lying below no prime in $S$ and such that $h(1/p)\ne0$
for the polynomial in Theorem~\ref{sys:bounded-extension}.
There are infinitely many eligible primes and only finitely
many roots of $h$. The theorem admits $1/p$ over some finite
extension $L/k$. At every prime of $L$ above $p$ this element
has negative valuation, and none of those primes belongs to
$S_L$. Thus $1/p\notin\mathcal O_{L,S_L}$.
\end{proof}

These corollaries concern a fixed formula whose witness field
is enlarged. They leave open the separate possibility of a
different defining formula for each number field.

\section{A finite collection of formulas}

\begin{theorem}\label{thm:finite-menu}
Let $\Phi_1,\ldots,\Phi_r$ be a finite nonempty list of
existential formulas over a number field $k$. Suppose
$\Phi_i(k)$ is infinite for every $i$.
There exists $B\ge1$ such that, for all but finitely many
$q\in\Q\setminus\Z$, there is a finite Galois extension
$L_q/k$ of degree at most $B$ in which every $\Phi_i(q)$ holds.
Thus none of the formulas defines $\mathcal O_{L_q}$ in $L_q$.
One may further require $L_q$ to have no real embedding, at the
cost of at most doubling $B$.
\end{theorem}
\begin{proof}
Apply Theorem~\ref{sys:bounded-extension} to each formula, obtaining $d_i$
and $h_i$. Exclude the finite union of the rational zero sets
of the $h_i$. For every remaining rational noninteger $q$,
choose the witnessing fields $E_{i,q}$.
The normal closure of $E_{i,q}/k$ has degree at most $d_i!$,
by its permutation action on the conjugates of a primitive
element. Take $L_q$ to be the compositum of these normal closures.
Then $L_q/k$ is Galois of degree at most
\[
                              B=\prod_{i=1}^r d_i!.
\]
Each witness tuple remains a witness in $L_q$: polynomial
equalities and nonzero conditions are preserved by field
extensions. A rational algebraic integer is an integer, so
$q\notin\mathcal O_{L_q}$. Adjoining $i$ gives the last assertion.
\end{proof}

\begin{corollary}
A fixed finite collection of formulas over $k$, each correct for
$\mathcal O_k$ in $k$, cannot supply a definition of the ring of
integers in every finite extension by choosing one formula from
that collection for each extension.
\end{corollary}
\begin{proof}
Each formula has infinitely many inputs over $k$. The extension
in Theorem~\ref{thm:finite-menu} makes every formula in the
collection incorrect simultaneously.
\end{proof}

\section{A sharper extension bound for one polynomial}

\begin{proposition}\label{sys:hypersurface-extension-bound}
Let $k$ be an infinite field and
$P\in k[T,Y_1,\ldots,Y_m]$. Suppose that $P$ depends on the
witness coordinate $Y_j$, and put $\delta=\deg_{Y_j}P\ge1$.
There is a nonzero $h\in k[T]$ such that for every $a\in k$
with $h(a)\ne0$, the equation $P(a,\mathbf y)=0$ has a solution
over an extension of $k$ of degree at most $\delta$.
\end{proposition}
\begin{proof}
Regard $P$ as a polynomial in $Y_j$. Its leading coefficient
is a nonzero polynomial in $T$ and the remaining witness
coordinates. Since $k$ is infinite, choose constants for those
remaining coordinates so that the leading coefficient becomes
a nonzero polynomial $h(T)$. The specialization is a polynomial
$Q(T,Z)\in k[T,Z]$ with $\deg_ZQ=\delta$. For $h(a)\ne0$,
choose any irreducible factor of the degree-$\delta$ polynomial
$Q(a,Z)$. Adjoining a root gives an extension of degree at
most $\delta$ and a witness for the original equation.
\end{proof}

\begin{corollary}\label{sys:single-polynomial-extension}
Let $k$ be a number field, and suppose a single nonzero polynomial
$P(T,\mathbf Y)$ has an infinite input projection over $k$.
Then $P$ depends on at least one witness coordinate. In
Corollary~\ref{sys:no-uniform-integers}, the bound on the extension
degree can be taken to be
\[
            D=\min_{j:\,\deg_{Y_j}P>0}\deg_{Y_j}P
               \le \deg P.
\]
In particular an unchanged quartic equation cannot define the
rings of integers over every extension of degree at most four.
If it has an essential witness of individual degree two, it
already fails over an extension of degree at most two.
\end{corollary}
\begin{proof}
If $P$ were independent of all witness coordinates, its input
projection would be the finite zero set of a nonzero polynomial
in $T$. Choose an essential coordinate of least positive
individual degree, apply Proposition~\ref{sys:hypersurface-extension-bound},
and take any rational noninteger outside its finite exceptional
set. The same argument applies to all such rational nonintegers.
\end{proof}

When the polynomial already defines an integral subset over
$k$, the failure in the last sentence cannot occur over $k$
itself; a witnessing extension of degree at most two must then
have degree exactly two. This is an obstruction to transferring
the equation unchanged, rather than a counterexample to
field-dependent Diophantine definability.


For a finite collection of nonzero polynomials with infinite input
projections, set
\[
 d_i=\min_{j:\,\deg_{Y_j}P_i>0}\deg_{Y_j}P_i.
\]
Proposition~\ref{sys:hypersurface-extension-bound} supplies the extensions
used in Theorem~\ref{thm:finite-menu}, so its Galois degree bound may be
chosen as $B=\prod_i d_i!$.

\section{An explicit quadratic extension}

Let $P_{392}$ be the quartic specified by the complete coefficient
and quadratic-residual tables in~\cite{QuarticData}. Coordinates are numbered $1,\ldots,392$,
with input $X_1$. It is the sum of squares of $364$ integer quadratic
residuals. Put
\[
 A=80989^{54929},\qquad
 D=16P_{392}(\tfrac12,0,\ldots,0).
\]
This formula defines a fixed integer without requiring its decimal
expansion. The exact computation in the supplied coefficient data gives
\[
 D=191779334771151081400576\,A^2+D_0,\qquad D_0>0,
\]
where the complete decimal value of $D_0$ is recorded in
\nolinkurl{computation/quartic/final392/field-specialization392.json}.

\begin{proposition}\label{sys:explicit-quartic-extension}
Let $L=\mathbb Q(\theta)$ with $\theta^2=-D$. Set $X_1=1/2$,
$X_{293}=\theta/4$, and every other witness coordinate to zero.
Then $P_{392}=0$ over $L$, but $X_1$ is not an algebraic integer.
Thus this unchanged quartic does not define $\mathcal O_L$ in $L$.
\end{proposition}
\begin{proof}
Leave coordinate $293$ free, calling it $Y$, and substitute the
specified values in the quadratic residuals. Of the $364$ residuals,
$302$ vanish, $61$ become constants, and the residual with source
identifier $493$, the square-output definition, becomes $-Y$.
Every constant belongs to $\tfrac14\mathbb Z[A]$, and one has absolute value $1/4$.
These statements follow by inspection of the explicitly listed
integer terms; the separate checker also verifies them by exact
substitution. Squaring and adding therefore gives
\[
 P_{392}(\tfrac12\mathbf e_1+Y\mathbf e_{293})
                      =Y^2+D/16.
\]
In particular $D$ is a positive integer. The independent checker
obtains this same identity both from all quadratic residuals and from
the complete quartic coefficient table, including its fixed value of $A$.
Since $D>0$, the polynomial $Z^2+D$ is irreducible over $\mathbb Q$,
so $L$ is an imaginary quadratic field. Taking $Y=\theta/4$ makes
the displayed value zero. The rational number $1/2$ is not an
algebraic integer, as proved in Corollary~\ref{sys:no-uniform-integers}.
\end{proof}


The exact substitution is verified independently from the quadratic
residuals and from their expanded quartic in
\nolinkurl{computation/quartic/check_field_specialization392.py}.
It gives a counterexample to transferring this specified equation
to the displayed field. It gives no counterexample to the existence
of a different definition of that field's ring of integers.

\begin{thebibliography}{99}
\bibitem{StacksChevalley}
The Stacks Project Authors, \emph{The Stacks Project},
Chevalley's theorem, Tag 00FE.
\url{https://stacks.math.columbia.edu/tag/00FE}.

\bibitem{QuarticData}
\emph{Complete coefficient and quadratic-residual tables for $P_{392}$},
exact project data,
\nolinkurl{computation/quartic/final392/polynomial392.json.gz} and
\nolinkurl{computation/quartic/final392/quadratic-system392.json.gz}.
\end{thebibliography}
\end{document}
